\uuid{V8Nb}
\exo7id{3876}
\titre{exo7 3876}
\auteur{quercia}
\organisation{exo7}
\datecreate{2010-03-11}
\isIndication{false}
\isCorrection{false}
\chapitre{Continuité, limite et étude de fonctions réelles}
\sousChapitre{Etude de fonctions}
\cpge{mpsi-mp2i.limites-continuite.b, pcsi.limites-continuite.b, ptsi.limites-continuite.b, mpsi-mp2i.limites-continuite.a, pcsi.limites-continuite.a, ptsi.limites-continuite.a, mpsi-mp2i.limites-continuite.c, pcsi.limites-continuite.c, ptsi.limites-continuite.c, bcpst1.limites-continuite.c, bcpst1.limites-continuite.a, bcpst1.fonctions-reelles-usuelles.a}
\module{Analyse}
\niveau{L1}
\difficulte{}

\contenu{
\question{
On pose : 
$$\begin{cases} f(x) = |x|\left({2+\sin\left(\frac1x\right)}\right) &\text { si } x \ne 0,\cr
                   f(0) = 0.\cr\end{cases}$$

Montrer que $f$ est continue, minimale en $0$, mais pour tout $\varepsilon > 0$,
$f_{|[0,\varepsilon]}$ n'est pas monotone.
}
}
